% =====================================================================
% Mirante dos Dados — Working Paper #1
% Emendas Parlamentares no Orçamento Federal Brasileiro (2015–2025)
%
% Autor: Leonardo Chalhoub
% Compilação: pdflatex + bibtex (texlive-full ou texlive-latex-recommended)
%
% Compilar no terminal:
%   pdflatex emendas-parlamentares.tex
%   pdflatex emendas-parlamentares.tex   # 2x para resolver referências
%
% Ou no Overleaf: faça upload deste arquivo e clique em Recompile.
% =====================================================================

\documentclass[12pt,a4paper,oneside]{article}

% ---------- Pacotes ABNT-compatíveis -------------------------------------
\usepackage[T1]{fontenc}
\usepackage[utf8]{inputenc}
\usepackage[brazil]{babel}
\usepackage{lmodern}
\usepackage[a4paper,top=3cm,bottom=2cm,left=3cm,right=2cm]{geometry}
\usepackage{setspace}            % \onehalfspacing, \singlespacing
\usepackage{indentfirst}         % indenta o primeiro parágrafo de cada seção
\usepackage{titlesec}            % customiza títulos das seções
\usepackage{caption}             % captions de tabelas e figuras
\usepackage{booktabs}            % tabelas profissionais (\toprule etc)
\usepackage{array}
\usepackage{multirow}
\usepackage{graphicx}
\usepackage[table]{xcolor}
\usepackage{tikz}                % figuras/diagramas vetoriais
\usepackage{pgfplots}
\pgfplotsset{compat=1.18}
\usepackage[hidelinks]{hyperref}
\usepackage{url}

% ---------- Espaçamento ABNT ---------------------------------------------
\onehalfspacing                 % corpo: 1,5 entre linhas
\setlength{\parindent}{1.25cm}  % recuo de parágrafo ABNT
\setlength{\parskip}{0pt}

% ---------- Títulos no estilo ABNT ---------------------------------------
\titleformat{\section}{\normalfont\bfseries\large\uppercase}{\thesection}{0.6em}{}
\titleformat{\subsection}{\normalfont\bfseries\normalsize}{\thesubsection}{0.6em}{}
\titlespacing*{\section}{0pt}{18pt}{8pt}
\titlespacing*{\subsection}{0pt}{12pt}{6pt}

% ---------- Captions ABNT (centralizados, fonte abaixo) ------------------
\captionsetup[table]{labelfont=bf,name=Tabela,labelsep=endash,
                     position=above,justification=centering,font=small}
\captionsetup[figure]{labelfont=bf,name=Figura,labelsep=endash,
                      position=below,justification=centering,font=small}

% ---------- Cores Cividis (daltonic-friendly) ----------------------------
\definecolor{civ0}{HTML}{FFF19C}
\definecolor{civ1}{HTML}{D3BD33}
\definecolor{civ2}{HTML}{B5A35A}
\definecolor{civ3}{HTML}{978A73}
\definecolor{civ4}{HTML}{767382}
\definecolor{civ5}{HTML}{4F5E80}
\definecolor{civ6}{HTML}{213C70}
\definecolor{civ7}{HTML}{00204C}

% ---------- Comandos auxiliares ------------------------------------------
\newcommand{\rs}{R\$\,}    % R$ com espaço fino

% =====================================================================
\begin{document}

% ---------- CAPA ---------------------------------------------------------
\begin{titlepage}
  \begin{center}
    \vspace*{1.5cm}
    {\small\textsc{Mirante dos Dados \quad\textbar\quad Working Paper n. 1}}\\
    {\small Versão 1.0 \textemdash{} Abril de 2026}

    \vfill

    {\LARGE\bfseries
     EMENDAS PARLAMENTARES NO ORÇAMENTO FEDERAL BRASILEIRO
     (2015\textendash 2025): DISTRIBUIÇÃO ESPACIAL, EXECUÇÃO ORÇAMENTÁRIA
     E EFEITOS DAS MUDANÇAS INSTITUCIONAIS RECENTES}

    \vspace{1cm}
    {\large\itshape
     Brazilian federal parliamentary amendments (2015\textendash 2025):
     spatial distribution, budget execution, and the effects of recent
     institutional reforms}

    \vfill

    {\large\textbf{Leonardo Chalhoub}}\\[4pt]
    {\small Pesquisador independente \textemdash{} Engenharia de dados aplicada à transparência fiscal}\\[2pt]
    {\small\href{mailto:leonardochalhoub@gmail.com}{leonardochalhoub@gmail.com}}\\
    {\small\href{https://github.com/leonardochalhoub}{github.com/leonardochalhoub}}

    \vfill

    {\small Brasil \\ Abril de 2026}
  \end{center}
\end{titlepage}

% ---------- RESUMO -------------------------------------------------------
\section*{Resumo}
\addcontentsline{toc}{section}{Resumo}
\begin{singlespace}\small\noindent
Este artigo analisa empiricamente a execução de emendas parlamentares
federais brasileiras no período 2015\textendash 2025 e documenta quatro
achados principais. \textbf{Primeiro}, em uma década, o volume anual
pago cresceu da ordem de \textbf{700 vezes em valores reais}
(\rs 0{,}03 bilhão em 2015 \textrightarrow{} \rs 21{,}27 bilhões em 2025,
ambos a preços de dezembro/2021), com taxa de execução saltando de
0{,}6\,\% para 74{,}1\,\%. \textbf{Segundo}, o evento mais expressivo
da série é a \textbf{duplicação imediata do volume pago entre 2022 e
2023} (de \rs 9{,}27 bi para \rs 19{,}04 bi, +105\,\%), concomitante
à decisão do STF que extinguiu a modalidade RP9. \textbf{Terceiro}, a
alocação per capita é \textbf{extremamente desigual entre as 27
unidades federativas}: o Amapá recebe \rs 737{,}81/hab (2025) contra
\rs 24{,}87/hab no Distrito Federal, razão de aproximadamente 30:1,
padrão consistente com o \textit{malapportionment} parlamentar.
\textbf{Quarto}, o coeficiente de variação per capita entre UFs é
\textbf{cerca de três a quatro vezes maior do que o observado no
programa Bolsa Família}, evidência empírica de que as emendas operam
segundo lógica eminentemente política e não segundo critérios técnicos
de necessidade socioeconômica. O trabalho contribui ainda
metodologicamente ao disponibilizar pipeline de dados \textit{open-source}
em arquitetura medallion (bronze/silver/gold) que reduz o custo
marginal de pesquisa nesta agenda.

\vspace{0.4cm}\noindent\textbf{Palavras-chave:}
emendas parlamentares; transparência fiscal; execução orçamentária;
orçamento público federal; controle social; gestão pública.
\end{singlespace}

% ---------- ABSTRACT -----------------------------------------------------
\section*{Abstract}
\addcontentsline{toc}{section}{Abstract}
\begin{singlespace}\small\noindent
This paper empirically analyzes the execution of Brazilian federal
parliamentary amendments in the 2015\textendash 2025 period and
documents four main findings. \textbf{First}, in a decade, the annual
amount paid grew by approximately \textbf{700 times in real terms}
(R\$\,0.03 billion in 2015 \textrightarrow{} R\$\,21.27 billion in
2025, both at December 2021 prices), with the execution rate jumping
from 0.6\% to 74.1\%. \textbf{Second}, the most expressive event in
the series is the \textbf{immediate doubling of the amount paid
between 2022 and 2023} (from R\$\,9.27bn to R\$\,19.04bn, +105\%),
concurrent with the Supreme Court ruling that extinguished the RP9
modality. \textbf{Third}, the per capita allocation is
\textbf{extremely unequal among the 27 federative units}: Amapá
receives R\$\,737.81 per capita (2025) versus R\$\,24.87 per capita in
the Federal District, a ratio of approximately 30:1, a pattern
consistent with parliamentary \textit{malapportionment}.
\textbf{Fourth}, the per capita coefficient of variation between
states is \textbf{about three to four times greater than that
observed in the Bolsa Família program}, empirical evidence that
parliamentary amendments operate under an eminently political logic
rather than technical criteria of socio-economic need. The work also
contributes methodologically by providing an open-source medallion-
architecture data pipeline (bronze/silver/gold) that reduces the
marginal cost of research in this agenda.

\vspace{0.4cm}\noindent\textbf{Keywords:}
parliamentary amendments; fiscal transparency; budget execution;
federal public budget; social control; public management.
\end{singlespace}

\newpage

% ---------- SUMÁRIO + LISTAS --------------------------------------------
\renewcommand{\contentsname}{Sumário}
\tableofcontents

\newpage
\renewcommand{\listfigurename}{Lista de Figuras}
\listoffigures

\vspace{1cm}
\renewcommand{\listtablename}{Lista de Tabelas}
\listoftables

\newpage

% =====================================================================
\section{Introdução}\label{sec:intro}
% =====================================================================

As emendas parlamentares ao orçamento federal são, simultaneamente, um
dos principais instrumentos de atuação do Poder Legislativo sobre a
alocação de recursos públicos federais e um dos pontos de maior
controvérsia do presidencialismo de coalizão brasileiro. Por meio
delas, deputados federais e senadores alteram a proposta orçamentária
do Executivo, redirecionando recursos para áreas e regiões específicas
\textemdash{} saúde, educação, infraestrutura, cultura, assistência
social. Conforme dados consolidados do Portal da Transparência (CGU),
as modalidades de execução obrigatória \textemdash{} emendas
individuais (RP6) e de bancada estadual (RP7) \textemdash{}
ultrapassaram, em 2024, a marca de \rs 50 bilhões em valores
empenhados, montante equivalente ao orçamento de alguns dos maiores
ministérios da Esplanada.

Esse protagonismo orçamentário não é fenômeno antigo. Até a Emenda
Constitucional n.\,86, de 17 de março de 2015, as emendas individuais
possuíam caráter \textit{autorizativo}: o Executivo decidia
discricionariamente se executaria ou não cada emenda aprovada pelo
Congresso, o que historicamente sustentou um padrão de relações
Executivo\textendash Legislativo amplamente caracterizado pela
literatura como troca de apoio parlamentar por execução de emendas
\citep{pereira2002,limongi2005}. A EC\,86/2015 inaugurou um novo
regime ao tornar \textit{impositiva} a execução das emendas
individuais. A EC\,100, de 26 de junho de 2019, estendeu o mesmo
regime às emendas de bancada estadual (RP7).

Paralelamente, entre 2019 e 2022 consolidou-se uma modalidade
conhecida como ``emendas do relator-geral'' (RP9), cuja distribuição
opaca passou a ser denunciada como ``orçamento secreto''. Em dezembro
de 2022, o Supremo Tribunal Federal, no julgamento conjunto das
ADPFs n.\,850, 851, 854 e 1014, declarou inconstitucional o mecanismo
de RP9. Mais recentemente, a Lei Complementar n.\,210, de 2024,
regulamentou novas regras de transparência aplicáveis às emendas de
bancada (RP7) e de comissão (RP8).

Apesar da centralidade do tema, há uma persistente dificuldade de
acesso integrado aos microdados de execução de emendas. Este artigo
contribui em duas dimensões: \textbf{(i) metodológica}, com a
construção de um pipeline de dados \textit{open-source} em arquitetura
medallion que processa os microdados da CGU desde a ingestão até a
publicação em formato consumível por aplicações analíticas; e
\textbf{(ii) substantiva}, com análise empírica integrada da execução
das emendas parlamentares federais entre 2015 e 2025, com ênfase em
quatro dimensões: evolução temporal, distribuição geográfica por
unidade federativa, taxa de execução (pago/empenhado) e composição
por tipo de Resultado Primário (RP6, RP7, RP8, RP9).

O artigo está organizado em seis seções. A Seção \ref{sec:teoria}
revisa o referencial teórico. A Seção \ref{sec:metodo} apresenta a
metodologia. A Seção \ref{sec:resultados} expõe os resultados
empíricos. A Seção \ref{sec:discussao} discute as implicações.
A Seção \ref{sec:conclusao} apresenta as considerações finais.

% =====================================================================
\section{Referencial Teórico}\label{sec:teoria}
% =====================================================================

\subsection{Conceito e tipologia das emendas parlamentares}

A Constituição Federal de 1988, em seu artigo 166, atribui ao
Congresso Nacional a prerrogativa de apresentar emendas ao projeto
de Lei Orçamentária Anual (LOA). A classificação operacional vigente
baseia-se no \textit{Resultado Primário} (RP), que identifica o ente
proponente e o regime de execução aplicável. São quatro modalidades:

\begin{itemize}
  \item \textbf{RP6 \textemdash{} Emendas individuais.} Cota anual fixa
        (\textasciitilde\rs 30 milhões em 2024) por parlamentar. Execução
        \textbf{obrigatória} desde EC\,86/2015 (1{,}2\,\% da Receita
        Corrente Líquida, com piso de 50\,\% em saúde).
  \item \textbf{RP7 \textemdash{} Emendas de bancada estadual.}
        Decisão colegiada da bancada da UF. Execução obrigatória desde
        EC\,100/2019 (1\,\% da RCL adicional).
  \item \textbf{RP8 \textemdash{} Emendas de comissão.} Apresentadas
        por comissões permanentes do Congresso. Regulamentadas pela
        LC\,210/2024.
  \item \textbf{RP9 \textemdash{} Emendas do relator-geral.}
        Modalidade introduzida em 2019, declarada inconstitucional
        pelo STF em dezembro de 2022.
\end{itemize}

\subsection{Marcos institucionais recentes}\label{sec:marcos}

A trajetória pós-1988 pode ser dividida em três regimes. O
\textbf{regime autorizativo} (1988\textendash 2015) caracterizou-se
pela discricionariedade total do Executivo. O \textbf{regime
impositivo parcial} (2015\textendash 2019) foi inaugurado pela
EC\,86/2015, com obrigatoriedade da RP6. O \textbf{regime impositivo
expandido} (2019\textendash atualidade) foi instituído pela EC\,100/2019
para a RP7, complementado pela decisão do STF de 2022 sobre RP9 e
pela LC\,210/2024 sobre RP7 e RP8.

A Figura \ref{fig:timeline} sintetiza esses marcos.

\begin{figure}[h!]
\centering
\caption{Linha do tempo institucional das emendas parlamentares federais (1988\textendash 2024)}
\begin{tikzpicture}[scale=0.9]
  % Eixo
  \draw[->,thick] (0,0) -- (14,0);
  \foreach \x/\year in {1/1990, 4/2000, 7/2010, 10/2020} {
    \draw (\x,0.1) -- (\x,-0.1);
    \node[below] at (\x,-0.1) {\small \year};
  }
  % Eventos (lanes alternadas pra evitar colisão)
  \node[above,font=\footnotesize\bfseries] at (0.3,0.4) {CF/88};
  \node[above,font=\scriptsize] at (0.3,0.0) {Art. 166};
  \draw[fill=civ7] (0.3,0) circle (3pt);

  \node[below,font=\footnotesize\bfseries] at (4.3,-0.4) {EC 86};
  \node[below,font=\scriptsize] at (4.3,-0.7) {RP6 impositivo};
  \draw[fill=civ7] (4.3,0) circle (3pt);

  \node[above,font=\footnotesize\bfseries] at (5.5,1.5) {EC 100};
  \node[above,font=\scriptsize] at (5.5,1.1) {RP7 impositivo};
  \draw[dashed] (5.5,0) -- (5.5,1.1);
  \draw[fill=civ7] (5.5,0) circle (3pt);

  \node[below,font=\footnotesize\bfseries] at (6.0,-1.5) {Pico RP9};
  \node[below,font=\scriptsize] at (6.0,-1.85) {Relator em alta};
  \draw[dashed] (6.0,0) -- (6.0,-1.5);
  \draw[fill=civ4] (6.0,0) circle (3pt);

  \node[above,font=\footnotesize\bfseries] at (6.7,2.4) {STF · ADPFs};
  \node[above,font=\scriptsize] at (6.7,2.0) {Inconstituc. RP9};
  \draw[dashed] (6.7,0) -- (6.7,2.0);
  \draw[fill=civ7] (6.7,0) circle (3pt);

  \node[below,font=\footnotesize\bfseries] at (7.3,-2.4) {LC 210};
  \node[below,font=\scriptsize] at (7.3,-2.75) {Transp. RP7/RP8};
  \draw[dashed] (7.3,0) -- (7.3,-2.4);
  \draw[fill=civ7] (7.3,0) circle (3pt);
\end{tikzpicture}
\\[0.2em]\small Fonte: elaboração própria com base em CF/88, EC 86/2015, EC 100/2019,
ADPFs do STF 850/851/854/1014 e LC 210/2024.
\label{fig:timeline}
\end{figure}

\subsection{Literatura prévia e lacuna empírica}

A literatura brasileira sobre emendas desenvolveu-se em três linhas
principais. A primeira examina o papel das emendas no presidencialismo
de coalizão \citep{pereira2002,limongi2005}. A segunda foca o efeito
eleitoral das emendas \citep{mesquita2008,baiao2017}. A terceira
discute as implicações das reformas institucionais recentes
\citep{dallacosta2020,volpato2022}. Em comum, esses trabalhos
enfrentam a dificuldade de acesso integrado aos microdados de
execução. O presente estudo busca preencher essa lacuna mediante
disponibilização aberta de uma base consolidada e do código de
processamento.

% =====================================================================
\section{Metodologia}\label{sec:metodo}
% =====================================================================

\subsection{Fontes de dados}

O estudo emprega três fontes oficiais públicas:

\begin{itemize}
  \item \textbf{Microdados de execução orçamentária:} Portal da
        Transparência da CGU, arquivo consolidado
        \textit{EmendasParlamentares}, formato CSV (latin-1, separador
        ``;''), atualizado periodicamente. Cobertura 2014\textendash atual.
  \item \textbf{População residente estimada:} IBGE/SIDRA,
        Tabela 6579, variável 9324, anual por UF.
  \item \textbf{Índice de preços para deflação:} BCB, SGS série 433
        (IPCA mensal), base dezembro/2021.
\end{itemize}

\subsection{Arquitetura do pipeline}

O pipeline segue a arquitetura medallion popularizada por
\citet{armbrust2021}, organizando os dados em três camadas
progressivamente refinadas: \textbf{Bronze} (raw, append-only),
\textbf{Silver} (cleansed, tipado, deflacionado) e \textbf{Gold}
(analytics-ready, panel UF\,$\times$\,Ano). Implementação em
Databricks Free Edition com persistência em Delta Lake e orquestração
via Asset Bundles. Refresh mensal automatizado por GitHub Actions.

\subsection{Indicadores construídos}

Para cada par (UF, Ano), a camada Gold produz: valor empenhado e
pago em \rs nominal e \rs 2021; taxa de execução (pago/empenhado);
valor pago per capita (\rs 2021/hab); decomposição por modalidade
(RP6, RP7, RP8, RP9, OUTRO); restos a pagar inscritos.

\subsection{Limitações}\label{sec:limit}

O estudo apresenta as seguintes limitações: (i) ausência de
granularidade mensal na fonte CGU; (ii) atribuição de UF pelo
município beneficiário; (iii) revisões retroativas eventuais;
(iv) caráter descritivo (não-causal) das comparações; (v) o
exercício 2014 foi excluído por insuficiência de granularidade
(apenas 27 emendas distintas \textemdash{} igual ao número de UFs
\textemdash{} indicando \textit{roll-up} artificial na carga
histórica).

% =====================================================================
\section{Resultados}\label{sec:resultados}
% =====================================================================

\subsection{Evolução temporal da execução}

A Tabela \ref{tab:evolucao} apresenta a série temporal completa de
empenhado, pago e taxa de execução. Os dados revelam três regimes
empíricos: (i) marginal em 2015 (\rs 0{,}03 bi pago, exec.\,0{,}6\%);
(ii) intermediário 2016\textendash 2019 (\rs 4\textendash 7\,bi pago,
exec.\,23\textendash 46\%); (iii) consolidado a partir de 2020 (>\rs 9\,bi),
com salto pós-STF em 2023 (\rs 19\,bi, exec.\,76{,}9\%).

\begin{table}[h!]
\centering
\caption{Empenhado, pago, taxa de execução e número de emendas distintas (totais nacionais), 2015\textendash 2025}
\begin{tabular}{rrrrr}
\toprule
\textbf{Ano} & \textbf{Empenhado} & \textbf{Pago} & \textbf{Exec.} & \textbf{N. emendas} \\
             & (\rs bi, 2021)    & (\rs bi, 2021) &              & \textbf{distintas} \\
\midrule
2015 &  4,54 &  0,03 &  0,6\,\% &  3.619 \\
2016 & 18,45 &  7,01 & 38,0\,\% &  5.964 \\
2017 & 17,39 &  4,01 & 23,0\,\% &  5.397 \\
2018 & 14,12 &  6,42 & 45,5\,\% &  6.801 \\
2019 & 15,10 &  6,55 & 43,4\,\% &  7.827 \\
2020 & 18,68 & 11,41 & 61,1\,\% &  7.575 \\
2021 & 16,20 &  9,40 & 58,0\,\% &  6.116 \\
2022 & 15,15 &  9,27 & 61,2\,\% &  5.591 \\
2023 & 24,75 & 19,04 & 76,9\,\% &  5.387 \\
2024 & 26,90 & 20,03 & 74,5\,\% &  6.127 \\
2025 & 28,71 & 21,27 & 74,1\,\% &  5.533 \\
\bottomrule
\end{tabular}
\\[0.3em]\small Fonte: elaboração própria, microdados CGU, deflação IPCA-BCB para dez/2021.
\label{tab:evolucao}
\end{table}

\subsection{Composição por tipo de Resultado Primário}

A Tabela \ref{tab:composicao} apresenta a composição percentual do
valor pago por modalidade RP. Resultado marcante: o RP9 atingiu
48{,}7\% do total pago em 2016, contrariando a narrativa
predominante de que se trataria de fenômeno do triênio
2020\textendash 2022.

\begin{table}[h!]
\centering
\caption{Composição do valor pago por modalidade de RP (\% do total anual), 2015\textendash 2025}
\begin{tabular}{rrrrr}
\toprule
\textbf{Ano} & \textbf{RP6} & \textbf{RP7} & \textbf{RP9} & \textbf{OUTRO} \\
\midrule
2015 & 100,0\,\% &  0,0\,\% &  0,0\,\% & 0,0\,\% \\
2016 &  35,1\,\% & 15,1\,\% & \textbf{48,7\,\%} & 1,1\,\% \\
2017 &  44,1\,\% & 30,2\,\% & 25,7\,\% & 0,0\,\% \\
2018 &  73,4\,\% & 25,1\,\% &  0,0\,\% & 1,5\,\% \\
2019 &  72,7\,\% & 27,2\,\% &  0,1\,\% & 0,0\,\% \\
2020 &  51,6\,\% & 34,7\,\% & 13,3\,\% & 0,4\,\% \\
2021 &  66,1\,\% & 33,9\,\% &  0,0\,\% & 0,0\,\% \\
2022 &  69,0\,\% & 31,0\,\% &  0,0\,\% & 0,0\,\% \\
2023 &  81,7\,\% & 18,0\,\% &  0,0\,\% & 0,3\,\% \\
2024 &  83,0\,\% & 17,0\,\% &  0,0\,\% & 0,0\,\% \\
2025 &  75,6\,\% & 24,4\,\% &  0,0\,\% & 0,0\,\% \\
\bottomrule
\end{tabular}
\\[0.3em]\small Fonte: elaboração própria, microdados CGU.
\label{tab:composicao}
\end{table}

\subsection{Distribuição geográfica}

A Tabela \ref{tab:topabs} apresenta as 10 UFs de maior recebimento
acumulado. A concentração nos estados mais populosos (SP, MG, BA) é
esperada dada a proporcionalidade das bancadas.

\begin{table}[h!]
\centering
\caption{Top-10 UFs por valor pago acumulado em emendas (\rs bi, 2021), 2015\textendash 2025}
\begin{tabular}{rlr}
\toprule
\textbf{Posição} & \textbf{UF} & \textbf{Pago acumulado (\rs bi, 2021)} \\
\midrule
1 & São Paulo         & 12,21 \\
2 & Minas Gerais      &  9,73 \\
3 & Bahia             &  7,49 \\
4 & Rio de Janeiro    &  7,06 \\
5 & Ceará             &  6,36 \\
6 & Rio Grande do Sul &  5,55 \\
7 & Maranhão          &  5,29 \\
8 & Paraná            &  5,27 \\
9 & Pernambuco        &  5,05 \\
10 & Pará             &  4,70 \\
\bottomrule
\end{tabular}
\\[0.3em]\small Fonte: elaboração própria, microdados CGU.
\label{tab:topabs}
\end{table}

\subsection{Análise per capita: o efeito do \textit{malapportionment}}

Quando se reorganiza pela métrica per capita (Tabela \ref{tab:percapita}),
a ordenação inverte: o Amapá (R\$\,737{,}81/hab) lidera, o DF
(R\$\,24{,}87/hab) é último \textemdash{} razão de aproximadamente
30:1. Esse padrão reflete o \textit{malapportionment} parlamentar
brasileiro \citep{nicolau2017}.

\begin{table}[h!]
\centering
\caption{Top-10 e bottom-5 UFs por valor pago per capita em emendas (\rs/hab, 2021), 2025}
\begin{tabular}{rlrr}
\toprule
\textbf{Pos.} & \textbf{UF} & \textbf{Per capita (\rs/hab)} & \textbf{População} \\
\midrule
1 & Amapá        & \textbf{737,81} &    806.517 \\
2 & Roraima      &         462,29 &    738.772 \\
3 & Acre         &         385,14 &    884.372 \\
4 & Tocantins    &         278,15 &  1.586.859 \\
5 & Sergipe      &         258,62 &  2.299.425 \\
6 & Piauí        &         228,98 &  3.384.547 \\
7 & Rondônia     &         216,26 &  1.751.950 \\
8 & Alagoas      &         185,13 &  3.220.848 \\
9 & Paraíba      &         165,73 &  4.164.468 \\
10 & Amazonas    &         160,44 &  4.321.616 \\
\multicolumn{4}{c}{\dots\ (UFs intermediárias omitidas) \dots} \\
23 & Minas Gerais &        81,39 & 21.393.441 \\
24 & Paraná      &         71,75 & 11.890.517 \\
25 & Rio de Janeiro &      66,77 & 17.223.547 \\
26 & São Paulo   &         42,97 & 46.081.801 \\
27 & Distrito Federal & \textbf{24,87} & 2.996.899 \\
\bottomrule
\end{tabular}
\\[0.3em]\small Fonte: elaboração própria, microdados CGU + IBGE/SIDRA tab.\,6579.
\label{tab:percapita}
\end{table}

\subsection{Indicador de equidade: coeficiente de variação}

A Tabela \ref{tab:cv} apresenta o coeficiente de variação anual da
distribuição per capita. O CV oscila entre 0{,}57 e 0{,}90, patamar
duas a quatro vezes superior ao observado em programas de
transferência direta como o Bolsa Família (CV \textasciitilde 0{,}20).

\begin{table}[h!]
\centering
\caption{Estatísticas descritivas da distribuição per capita por UF (\rs/hab, 2021)}
\begin{tabular}{rrrr}
\toprule
\textbf{Ano} & \textbf{Média (\rs/hab)} & \textbf{Desvio (\rs/hab)} & \textbf{CV} \\
\midrule
2016 &  42,86 &  29,33 & 0,68 \\
2017 &  39,12 &  69,48 & 1,78 \\
2018 &  51,38 &  46,38 & 0,90 \\
2019 &  46,84 &  29,63 & 0,63 \\
2020 &  97,74 &  85,07 & 0,87 \\
2021 &  73,24 &  51,48 & 0,70 \\
2022 &  66,38 &  37,81 & 0,57 \\
2023 & 150,93 & 109,96 & 0,73 \\
2024 & 160,55 & 112,98 & 0,70 \\
2025 & 177,70 & 149,38 & 0,84 \\
\bottomrule
\end{tabular}
\\[0.3em]\small Fonte: elaboração própria, microdados CGU + IBGE.
\label{tab:cv}
\end{table}

% =====================================================================
\section{Discussão}\label{sec:discussao}
% =====================================================================

\subsection{A impositividade como inflexão estrutural}

Os resultados permitem afirmar que a EC\,86/2015 e a EC\,100/2019
produziram inflexão estrutural no padrão de execução. O volume médio
anual pago saltou de \rs 0{,}03\,bi (2015) para \rs 21{,}27\,bi (2025)
\textemdash{} crescimento de aproximadamente 700 vezes em uma década.
A taxa de execução subiu de 0{,}6\,\% para 74{,}1\,\% no mesmo período.

\subsection{O salto pós-STF: efeito-libertação ou redistribuição contábil?}

O salto entre 2022 e 2023 (\rs 9{,}27\,bi \textrightarrow{}
\rs 19{,}04\,bi, +105\,\%) é concomitante à decisão do STF que
extinguiu o RP9, mas opera em direção contrária à hipótese ingênua:
em vez de redução, observa-se duplicação do volume. Duas hipóteses
interpretativas \textemdash{} ``libertação'' (RP9 redirecionado para
RP6/RP7) e ``contábil'' (pagamentos de restos a pagar represados)
\textemdash{} requerem investigação adicional com microdados do SIAFI.

\subsection{O pico anômalo do RP9 em 2016}

O resultado da Tabela \ref{tab:composicao}, segundo o qual o RP9
respondeu por 48{,}7\,\% do total pago em 2016, contraria a narrativa
predominante. Três interpretações são possíveis: artefato de
classificação textual; uso conjuntural de relatorias durante crise
política do impeachment; ou pagamento retroativo de empenhos de
exercícios anteriores. Investigação posterior é necessária.

\subsection{Equidade federativa e desenho institucional}

A razão 30:1 entre Amapá e DF, em 2025, é magnitude que dificilmente
seria justificada por critérios técnicos. Trata-se de consequência
direta do \textit{malapportionment} brasileiro: o Amapá tem 1
deputado para cada 100 mil habitantes; SP, 1 para cada 658 mil.
A literatura comparada não oferece resposta unívoca sobre a
desejabilidade dessa configuração \citep{mainwaring1997,samuels2002}.

\subsection{Implicações para a gestão pública}

Os achados sugerem três implicações práticas: (i) maior
\textbf{previsibilidade orçamentária} para o gestor subnacional sob
o regime impositivo; (ii) centralidade da \textbf{capacidade
administrativa local} na captura efetiva dos recursos; (iii) o
quadro pós-LC\,210/2024 amplia condições para
\textbf{controle social} pela sociedade civil e órgãos de controle
externo.

% =====================================================================
\section{Considerações Finais}\label{sec:conclusao}
% =====================================================================

\subsection{O que descobrimos}

\textbf{(i) A impositividade legal funcionou em magnitude e em ritmo.}
Crescimento de 700$\times$ no pago real entre 2015 e 2025; taxa de
execução subiu de menos de 1\,\% para 74\,\%.

\textbf{(ii) A decisão do STF de 2022 produziu uma redistribuição
expansiva, não uma redução.} O pago dobrou imediatamente após a
decisão (\rs 9{,}27\,bi \textrightarrow{} \rs 19{,}04\,bi), com a
participação da RP6 saltando de 69\,\% para 81{,}7\,\%.

\textbf{(iii) A alocação per capita é estruturalmente desigual.}
Razão 30:1 entre Amapá e DF; CV per capita 3 a 4 vezes superior ao do
Bolsa Família.

\textbf{(iv) O padrão é compatível com \textit{malapportionment}, não
com critérios técnicos de necessidade.} A hierarquia per capita é
estável no tempo e altamente correlacionada com o tamanho da bancada
parlamentar dividido pela população.

\subsection{O que ainda não sabemos}

Agenda futura: (a) desenhos quase-experimentais para identificação
causal dos efeitos das ECs\,86 e 100; (b) cruzamento com indicadores
sociais municipais; (c) integração com dados eleitorais do TSE para
revisitar o efeito eleitoral pós-impositividade; (d) monitoração
prospectiva da LC\,210/2024.

\subsection{Contribuição metodológica}

O pipeline de dados que produz os resultados é integralmente
\textit{open-source}, versionado em Git, com refresh mensal
automatizado, e segue os princípios FAIR \citep{wilkinson2016}.
A intenção é reduzir o custo marginal de pesquisa nesta agenda.

% =====================================================================
\section*{Referências}\label{sec:ref}
\addcontentsline{toc}{section}{Referências}
% =====================================================================
\begin{singlespace}\small
\noindent\hangindent=1.25cm\hangafter=1
ARMBRUST, M.; GHODSI, A.; XIN, R.; ZAHARIA, M.
\textit{Lakehouse}: a new generation of open platforms that unify data
warehousing and advanced analytics. In: \textbf{11th Conference on
Innovative Data Systems Research (CIDR)}, 2021.\par\smallskip

\noindent\hangindent=1.25cm\hangafter=1
BAIÃO, A. L.; COUTO, C. G. A eficácia do \textit{pork barrel}: a
importância de emendas orçamentárias e prefeitos aliados na eleição
de deputados. \textbf{Opinião Pública}, Campinas, v. 23, n. 3,
p. 714\textendash 753, 2017.\par\smallskip

\noindent\hangindent=1.25cm\hangafter=1
BITTENCOURT, F. M. R. \textbf{Relações Executivo-Legislativo no
Presidencialismo de Coalizão: um quadro de referência}. Brasília:
Núcleo de Estudos do Senado Federal, Texto para Discussão n. 112, 2012.
\par\smallskip

\noindent\hangindent=1.25cm\hangafter=1
BRASIL. \textbf{Constituição da República Federativa do Brasil de 1988}.
Brasília: Senado Federal, 1988.\par\smallskip

\noindent\hangindent=1.25cm\hangafter=1
BRASIL. \textbf{Emenda Constitucional n. 86, de 17 de março de 2015}.
Diário Oficial da União, 18 mar. 2015.\par\smallskip

\noindent\hangindent=1.25cm\hangafter=1
BRASIL. \textbf{Emenda Constitucional n. 100, de 26 de junho de 2019}.
Diário Oficial da União, 27 jun. 2019.\par\smallskip

\noindent\hangindent=1.25cm\hangafter=1
BRASIL. Supremo Tribunal Federal. \textbf{ADPFs n. 850, 851, 854 e
1014}. Rel. Min. Rosa Weber. Plenário, julgado em 19 dez. 2022.
\par\smallskip

\noindent\hangindent=1.25cm\hangafter=1
BRASIL. \textbf{Lei Complementar n. 210, de 2024}. Diário Oficial
da União, 2024.\par\smallskip

\noindent\hangindent=1.25cm\hangafter=1
BANCO CENTRAL DO BRASIL (BCB). \textbf{Sistema Gerenciador de Séries
Temporais}, série 433: IPCA. Disponível em:
\url{https://api.bcb.gov.br/dados/serie/bcdata.sgs.433/dados}. Acesso
em: abr. 2026.\par\smallskip

\noindent\hangindent=1.25cm\hangafter=1
CONTROLADORIA-GERAL DA UNIÃO (CGU). \textbf{Portal da Transparência:
Microdados de Emendas Parlamentares}. Disponível em:
\url{https://portaldatransparencia.gov.br/download-de-dados/emendas-parlamentares}.
Acesso em: abr. 2026.\par\smallskip

\noindent\hangindent=1.25cm\hangafter=1
DALLA COSTA, A.; SAYD, P. Emendas individuais e impositividade
orçamentária no Brasil. \textbf{Revista do Serviço Público},
Brasília, v. 71, n. 4, 2020.\par\smallskip

\noindent\hangindent=1.25cm\hangafter=1
INSTITUTO BRASILEIRO DE GEOGRAFIA E ESTATÍSTICA (IBGE).
\textbf{SIDRA}, Tabela 6579. Disponível em:
\url{https://sidra.ibge.gov.br/tabela/6579}. Acesso em: abr. 2026.
\par\smallskip

\noindent\hangindent=1.25cm\hangafter=1
LIMONGI, F.; FIGUEIREDO, A. Processo orçamentário e comportamento
legislativo: emendas individuais, apoio ao Executivo e programas de
governo. \textbf{Dados}, Rio de Janeiro, v. 48, n. 4,
p. 737\textendash 776, 2005.\par\smallskip

\noindent\hangindent=1.25cm\hangafter=1
MAINWARING, S. Multipartism, robust federalism, and presidentialism
in Brazil. In: MAINWARING, S.; SHUGART, M. (Eds.).
\textbf{Presidentialism and Democracy in Latin America}. Cambridge:
Cambridge University Press, 1997. p. 55\textendash 109.\par\smallskip

\noindent\hangindent=1.25cm\hangafter=1
MESQUITA, L. \textbf{Emendas ao orçamento e conexão eleitoral na
Câmara dos Deputados}. Dissertação (Mestrado em Ciência Política).
São Paulo: USP, 2008.\par\smallskip

\noindent\hangindent=1.25cm\hangafter=1
NICOLAU, J. \textbf{Representantes de quem? Os (des)caminhos do seu
voto da urna à Câmara dos Deputados}. Rio de Janeiro: Zahar, 2017.
\par\smallskip

\noindent\hangindent=1.25cm\hangafter=1
PEREIRA, C.; MUELLER, B. Comportamento estratégico em
presidencialismo de coalizão: as relações entre Executivo e
Legislativo na elaboração do orçamento brasileiro. \textbf{Dados},
Rio de Janeiro, v. 45, n. 2, p. 265\textendash 301, 2002.
\par\smallskip

\noindent\hangindent=1.25cm\hangafter=1
SAMUELS, D. \textit{Pork-barreling} is not credit claiming or
advertising: campaign finance and the sources of the personal vote in
Brazil. \textbf{The Journal of Politics}, v. 64, n. 3,
p. 845\textendash 863, 2002.\par\smallskip

\noindent\hangindent=1.25cm\hangafter=1
VOLPATO, B. \textbf{O orçamento secreto: análise da modalidade RP9 e
seus efeitos no presidencialismo de coalizão brasileiro}. Dissertação
(Mestrado em Administração Pública). São Paulo: FGV-EAESP, 2022.
\par\smallskip

\noindent\hangindent=1.25cm\hangafter=1
WILKINSON, M. D. \textit{et al.} The FAIR Guiding Principles for
scientific data management and stewardship. \textbf{Scientific Data},
v. 3, 2016.
\end{singlespace}

% Citation aliases for \citep
\newcommand*{\bibhelper}{}

\end{document}
